%
% Halaman Abstrak
%
% @author  Andreas Febrian
% @version 1.00
%

\chapter*{\normalsize ABSTRAK}
\begin{singlespace}
\vspace*{0.2cm}

\noindent\begin{tabular}{l l >{\raggedright\arraybackslash}p{10cm}}
Nama&: & \penulis \\ % Jangan diubah
Program Studi&: & \jenjang~\bidang \\ % Jangan diubah
Judul&: & \judul \\ % Jangan diubah
\end{tabular}

\vspace*{0.5cm}

\noindent

Wilayah strategis Banten, DKI Jakarta, dan Jawa Barat sebagai pusat ekonomi nasional mengalami perubahan struktural kompleks selama 2020--2024 akibat pandemi COVID-19 dan fase pemulihan yang tidak merata antar-kabupaten/kota; pendekatan multivariat konvensional memiliki keterbatasan dalam menangkap interaksi simultan antara dimensi spasial, indikator sosial-ekonomi, dan pergeseran temporal secara terpadu. Penelitian ini bertujuan mengevaluasi stabilitas longitudinal pada periode saat pandemi (2020--2021) dan pasca pandemi (2022--2024), serta memetakan trajektori spasio-temporal sebelas indikator pembangunan strategis yang meliputi aspek ekonomi (PDRB), kesejahteraan (Kemiskinan, Garis Kemiskinan, Pangan), pembangunan manusia (Pendidikan, Kesehatan), ketenagakerjaan, serta kesetaraan (Gender, Gini Ratio) di 41 kabupaten/kota pada ketiga provinsi tersebut. Metode STATIS (\textit{Structuration des Tableaux à Trois Indices de la Statistique} atau Strukturisasi Tabel dengan Tiga Indeks dalam Statistik) diterapkan melalui tahapan: (1)~standardisasi $Z$-\textit{score} terbobot demografis menggunakan jumlah penduduk sebagai bobot metrik, (2)~evaluasi inter-struktur menggunakan koefisien $RV$ untuk mengukur kemiripan antar-tabel tahunan, (3)~pembentukan struktur kompromi ($\mathbf{W}_c$) melalui pembobotan spektral optimal, (4)~validasi signifikansi via uji permutasi Monte Carlo, dan (5)~visualisasi lintasan (\textit{trajectory}) menggunakan \textit{Dynamic Biplot}; pengelompokan lanjutan diterapkan melalui algoritma \textit{K-Means} ($k=3$) pada koordinat kompromi yang direduksi via \textit{Principal Component Analysis} (PCA). Analisis inter-struktur menunjukkan konsistensi yang sangat tinggi dengan seluruh koefisien $RV$ melampaui ambang batas kritis dan terbukti signifikan secara statistik, membuktikan bahwa hierarki relatif antar-wilayah tetap stabil meskipun terjadi fluktuasi absolut akibat pandemi; struktur kompromi mampu menjelaskan sebagian besar total variansi melalui dua dimensi utama, yaitu Dimensi~1 yang merepresentasikan "Kapasitas Sosio-Ekonomi" (berkorelasi kuat dengan PDRB dan IPM) dan Dimensi~2 yang mencerminkan "Tekanan Pasar Tenaga Kerja" (didominasi Tingkat Pengangguran Terbuka). \textit{Dynamic Biplot} mengidentifikasi tiga pola trajektori: pemulihan berbentuk~V pada pusat metropolitan berbasis jasa, trajektori stagnan di kabupaten selatan Banten dan Jawa Barat, serta pertumbuhan progresif di wilayah penyangga; pengelompokan \textit{K-Means} menghasilkan tiga klaster yaitu "Maju," "Berkembang," dan "Tertinggal" yang memiliki keselarasan tinggi dengan kategori normatif IPM UNDP, mengindikasikan bahwa kesenjangan regional bersumber dari faktor struktural yang persisten. Penelitian ini dibatasi pada tiga provinsi di Jawa bagian Barat dengan sebelas indikator terpilih dan rentang waktu lima tahun (2020--2024); kebaruan penelitian ini terletak pada penerapan metode STATIS untuk analisis dinamika pembangunan regional pada periode saat dan pasca pandemi COVID-19 di konteks Jawa bagian Barat, yang diintegrasikan dengan sistem \textit{dashboard} web interaktif berbasis Python mencakup uji diagnostik, visualisasi klaster spasio-temporal, serta ekspor otomatis sebagai instrumen analitik bagi peneliti maupun pemangku kebijakan.

\bigskip

\noindent

\vspace*{0.5cm}

\noindent \bo{Kata kunci}:
	Analisis Data Tiga Arah, Data Longitudinal, \textit{K-Means Clustering}, PCA, Visualisasi Interaktif.
\end{singlespace}
\newpage 